%!TeX root=MemoriaTFG.tex

\chapter{Fonaments Teòrics}
\label{cap:marc}

Aquest capítol presenta els fonaments teòrics que sustenten el treball: els
conceptes bàsics de l'\ac{AR}, els algorismes emprats i l'estat de l'art en
jocs d'informació parcial. Cada secció s'orienta específicament als components
que apareixeran a la metodologia d'entrenament.

% ────────────────────────────────────────────────────────────────────────────
\section{Aprenentatge per Reforç}
\label{sec:ar}

L'\ac{AR} és el paradigma d'aprenentatge automàtic en el qual un agent aprèn
a actuar en un entorn mitjançant la interacció directa, guiat per senyals de
recompensa~\cite{sutton2018}. A diferència de l'aprenentatge supervisat, no
existeix cap conjunt d'exemples etiquetats ni un tercer que indiqui la
resposta correcta; a diferència de l'aprenentatge no supervisat, sí que hi ha
una senyal d'avaluació (la recompensa) que orienta l'aprenentatge.

Aquesta secció presenta els fonaments formals de l'\ac{AR}: la formalització
com a \ac{MDP}, les funcions de valor i les equacions de Bellman, el problema
d'exploració, la distinció entre algorismes \emph{on-policy} i
\emph{off-policy}, i l'extensió al cas d'observació parcial (\ac{POMDP}).

\subsection*{Procés de Decisió de Markov}

L'entorn es formalitza com un \acf{MDP}, definit per la tupla
$(S, A, P, R, \gamma)$ on cada component té el significat que recull la
Taula~\ref{tab:mdp-tuple}.

\begin{table}[H]
  \centering
  \begin{tabular}{ll}
    \toprule
    Component & Descripció \\
    \midrule
    $S$ & Conjunt d'estats de l'entorn \\
    $A$ & Conjunt d'accions disponibles per a l'agent \\
    $P(s' \mid s, a)$ & Probabilitat de transitar a $s'$ prenent l'acció $a$ des de $s$ \\
    $R(s, a, s')$ & Recompensa immediata esperada en la transició \\
    $\gamma \in [0,1)$ & Factor de descompte (recompenses futures vs.\ immediates) \\
    \bottomrule
  \end{tabular}
  \caption{Components de la tupla \acs{MDP}.}
  \label{tab:mdp-tuple}
\end{table}

Una \emph{política} $\pi$ és la regla de decisió de l'agent: en la forma
estocàstica, $\pi(a \mid s)$ expressa la probabilitat de seleccionar l'acció
$a$ quan l'estat és $s$; en la forma determinista, $\pi : S \to A$ assigna
directament una acció a cada estat. El cicle bàsic d'interacció és el
representat a la Figura~\ref{fig:cicle-ar}: a cada pas de temps $t$, l'agent
observa $s_t$, tria $a_t \sim \pi(\cdot \mid s_t)$, rep la recompensa $r_t$ i
l'entorn transita a $s_{t+1} \sim P(\cdot \mid s_t, a_t)$.

\begin{figure}[H]
  \centering
  \begin{tikzpicture}[
    bloc/.style={draw, rounded corners=4pt, minimum width=2.8cm,
                 minimum height=1.0cm, align=center, font=\small},
    >=Stealth, thick
  ]
    \node[bloc] (agent) {Agent\\[2pt]$\pi(a\mid s;\theta)$};
    \node[bloc, right=5cm of agent] (env) {Entorn};
    \draw[->] ([yshift=4pt]agent.east) --
              node[above, font=\small]{acció $a_t$}
              ([yshift=4pt]env.west);
    \draw[->] ([yshift=-4pt]env.west) --
              node[below, font=\small]{estat $s_{t+1}$,\quad recompensa $r_t$}
              ([yshift=-4pt]agent.east);
  \end{tikzpicture}
  \caption{Cicle d'interacció agent–entorn en \acs{AR}.}
  \label{fig:cicle-ar}
\end{figure}

L'objectiu de l'agent és trobar la política òptima $\pi^*$ que maximitzi el
\emph{retorn acumulat descomptat}:
\[
  G_t = \sum_{k=0}^{\infty} \gamma^k\, r_{t+k+1}.
\]
El factor $\gamma$ controla l'horitzó temporal efectiu: quan $\gamma \to 0$
l'agent és miop i prioritza les recompenses immediates; quan $\gamma \to 1$
valora el futur quasi tant com el present.

\FloatBarrier
\subsection*{Funcions de valor i equacions de Bellman}

Per saber si una política $\pi$ és bona, cal quantificar el retorn esperat
que l'agent obtindrà seguint-la. Per fer-ho s'utilitzen dues funcions de
valor complementàries: la \emph{funció de valor d'estat} $V^\pi(s)$ i la
\emph{funció de valor d'acció} $Q^\pi(s,a)$:
\[
  V^\pi(s) = \mathbb{E}_\pi\!\left[G_t \;\middle|\; s_t = s\right], \qquad
  Q^\pi(s,a) = \mathbb{E}_\pi\!\left[G_t \;\middle|\; s_t = s,\, a_t = a\right].
\]

Aquestes funcions satisfan les \emph{equacions de Bellman}, que les expressen
de forma recursiva a partir del pas immediatament següent:
\begin{align}
  V^\pi(s) &= \sum_{a} \pi(a \mid s)
              \sum_{s'} P(s' \mid s, a)
              \bigl[R(s,a,s') + \gamma\, V^\pi(s')\bigr], \label{eq:bellman-v}\\
  Q^\pi(s,a) &= \sum_{s'} P(s' \mid s, a)
               \bigl[R(s,a,s') + \gamma
               \sum_{a'}\pi(a' \mid s')\,Q^\pi(s',a')\bigr]. \label{eq:bellman-q}
\end{align}

Quan la política és l'òptima $\pi^*$, les equacions es simplifiquen
substituint la suma sobre accions per un màxim:
\begin{align}
  V^*(s)   &= \max_{a}\; \mathbb{E}\bigl[r + \gamma\, V^*(s')\mid s,a\bigr],
                         \label{eq:bellman-opt-v}\\
  Q^*(s,a) &= \mathbb{E}\bigl[r + \gamma \max_{a'} Q^*(s',a') \mid s,a\bigr].
                         \label{eq:bellman-opt-q}
\end{align}

La relació entre totes dues és $V^*(s) = \max_a Q^*(s,a)$, i la política
òptima s'obté directament com $\pi^*(s) = \arg\max_a Q^*(s,a)$. Quan l'espai
d'estats és massa gran per representar $Q^*$ en una taula, s'aproxima amb una
xarxa neuronal $Q(s,a;\theta)$, donant lloc a la família d'algorismes de
\ac{DRL}.

\FloatBarrier
\subsection*{Exploració i explotació}

Fins i tot amb les equacions de Bellman ben definides, aprendre $Q^*$ per
interacció directa planteja un dilema pràctic: l'agent ha de decidir si aplica
l'acció que creu millor (\textbf{explotació}) o en prova una de nova per
obtenir més informació (\textbf{exploració}). Sense exploració suficient,
l'agent queda atrapat en polítiques subòptimes; amb massa exploració, no
convergeix.

L'estratègia $\varepsilon$-\emph{greedy} és la solució més senzilla: amb
probabilitat $1-\varepsilon$ tria l'acció de valor màxim i amb probabilitat
$\varepsilon$ n'escull una uniformement a l'atzar. Durant l'entrenament,
$\varepsilon$ decau des d'un valor proper a 1 fins a un valor residual baix.
Els algorismes de gradient de política com \ac{PPO} adopten un enfocament
diferent: afegeixen un terme d'entropia $\mathcal{H}[\pi]$ a la funció
objectiu, penalitzant les polítiques massa deterministes i incentivant la
diversitat d'accions de forma contínua.

\FloatBarrier
\subsection*{Polítiques \emph{on-policy} i \emph{off-policy}}

Els algorismes \emph{on-policy} (com \ac{PPO}) aprenen únicament de dades
generades per la política actual; els \emph{off-policy} (com \ac{DQN}) poden
reutilitzar transicions passades emmagatzemades en un \emph{replay buffer},
cosa que millora l'eficiència de mostres però requereix correcció de la
distribució.

\FloatBarrier
\subsection*{Observació parcial: el model POMDP}

Els models \ac{MDP} vistos fins ara assumeixen que l'agent observa l'estat
complet del sistema en cada instant. En molts entorns reals, i en particular
al Truc, aquesta suposició no es compleix: els adversaris amaguen les seves
cartes i les intencions futures romanen ocultes.

Un \acf{POMDP} estén el \ac{MDP} amb dues components addicionals (un conjunt
d'observacions $\Omega$ i una funció d'emissió $O$), donant la tupla completa
$(S, A, P, R, \gamma, \Omega, O)$. La funció $O(o \mid s', a)$ especifica la
probabilitat de rebre l'observació $o \in \Omega$ després de transitar a
l'estat $s'$ prenent l'acció $a$. L'agent no veu $s_t$ directament, sinó una
observació parcial $o_t$ que en conté menys informació, tal com il·lustra la
Figura~\ref{fig:pomdp}.

\begin{figure}[H]
  \centering
  \begin{tikzpicture}[
    state/.style={draw, circle, minimum size=1.0cm, font=\small},
    obs/.style  ={draw, circle, minimum size=1.0cm, font=\small, dashed},
    act/.style  ={font=\small},
    >=Stealth, thick
  ]
    \foreach \i/\t in {1/$s_{t-1}$, 2/$s_t$, 3/$s_{t+1}$}{
      \node[state] (s\i) at ({(\i-1)*3.2}, 0) {\t};
    }
    \foreach \i/\t in {1/$o_{t-1}$, 2/$o_t$, 3/$o_{t+1}$}{
      \node[obs] (o\i) at ({(\i-1)*3.2}, -2) {\t};
    }
    \draw[->] (s1) -- (s2) node[midway, above, font=\scriptsize]{$a_{t-1}$};
    \draw[->] (s2) -- (s3) node[midway, above, font=\scriptsize]{$a_t$};
    \draw[->, dashed] (s1) -- (o1) node[midway, right, font=\scriptsize]{$O$};
    \draw[->, dashed] (s2) -- (o2) node[midway, right, font=\scriptsize]{$O$};
    \draw[->, dashed] (s3) -- (o3) node[midway, right, font=\scriptsize]{$O$};
  \end{tikzpicture}
  \caption{Model gràfic del \acs{POMDP}: els estats $s_t$ (cercles continus) són
           ocults; l'agent només observa $o_t$ (cercles discontinus) generats per
           la funció d'emissió $O$.}
  \label{fig:pomdp}
\end{figure}

La política òptima en un \ac{POMDP} ha de raonar sobre la
\emph{distribució de creença} (\emph{belief state}) $b_t(s) = P(s_t = s \mid
o_{1:t}, a_{1:t-1})$. Quan es rep una nova observació $o_{t+1}$, la creença
s'actualitza aplicant el teorema de Bayes:
\[
  b_{t+1}(s') \;\propto\; O(o_{t+1} \mid s', a_t)
  \sum_{s \in S} P(s' \mid s, a_t)\, b_t(s).
\]
En la pràctica, mantenir $b_t$ explícitament és intractable per a espais
d'estats grans. En el seu lloc s'utilitzen xarxes recurrents (\acs{LSTM},
\acs{GRU}) que comprimeixen l'historial d'observacions en un vector d'estat
ocult, actuant com a representació implícita de la creença.

L'entorn del Truc és precisament un \ac{POMDP}: l'agent rep una observació
que inclou la seva mà i l'historial d'apostes visible, però no les cartes dels
adversaris.

\FloatBarrier

% ─────────────────────────────────────────────────────────────────────────────
\section{Estat de l'art en jocs d'informació parcial}
\label{sec:estat-art}

El pòquer ha estat el banc de prova principal per a la IA en jocs d'informació
parcial, i els sistemes basats en Minimització del Penediment Contra-factual
(\emph{Counterfactual Regret Minimization}, CFR) han assolit resultats
sobrehumans.

\textbf{Libratus}~\cite{brown2018libratus} (Brown \& Sandholm, 2017) guanya
jugadors professionals de Texas Hold'em de dos jugadors. CFR minimitza
iterativament el penediment contra-factual (la millora hipotètica si s'hagués
pres una acció diferent) i convergeix a l'equilibri de Nash. Libratus resol
sub-jocs en temps real per refinar la seva estratègia durant la partida.

\textbf{Pluribus}~\cite{brown2019pluribus} (Brown \& Sandholm, 2019) estén
aquest resultat al pòquer de sis jugadors combinant Monte Carlo CFR, una
estratègia de base (\emph{blueprint}) obtinguda per self-play i la resolució
de sub-jocs en temps real. És la primera IA a superar humans en un entorn
multi-jugador d'informació parcial.

\subsection*{Limitacions dels enfocaments CFR}

Malgrat els seus èxits, CFR presenta tres limitacions que motiven l'ús de
\ac{DRL} en aquest treball. El cost creix exponencialment amb la profunditat
de l'arbre i el nombre de jugadors, cosa que exigeix abstraccions manuals
específiques per a cada joc. A més, aquestes abstraccions no es transfereixen:
redissenyar-les per al Truc (amb apostes menys estructurades i senyes) seria
tan costós com el problema original. Finalment, CFR no aprèn representacions
compartides, de manera que no es beneficia de la generalització de les xarxes
neuronals.

L'enfocament \ac{DRL} (DQN, PPO, NFSP) evita modelar l'arbre explícitament i
generalitza a nous jocs sense redissenyar abstraccions, a canvi d'una
convergència a Nash menys garantida teòricament.

% ─────────────────────────────────────────────────────────────────────────────
\section{Deep Q-Network}
\label{sec:dqn}

\acf{DQN}~\cite{mnih2015dqn} fou el primer algoritme de \ac{DRL} a demostrar
aprenentatge de nivell humà en entorns d'alta dimensionalitat (jocs Atari),
combinant Q-learning amb xarxes neuronals profundes.

En lloc d'una taula $Q(s,a)$, \ac{DQN} aproxima la funció de valor d'acció
amb una xarxa neuronal $Q(s,a;\theta)$ parametritzada per $\theta$. Aprèn
minimitzant la pèrdua:
\[
  \mathcal{L}(\theta) =
  \mathbb{E}_{(s,a,r,s') \sim \mathcal{D}}\!\left[
    \bigl(r + \gamma \max_{a'} Q(s',a';\theta^{-}) - Q(s,a;\theta)\bigr)^2
  \right].
\]

Dues innovacions clau estabilitzen l'entrenament:

\begin{itemize}
  \item \textbf{Experience Replay}: les transicions $(s,a,r,s')$ s'emmagatzemen
        en un buffer $\mathcal{D}$ i es mostregen aleatòriament. Això trenca la
        correlació temporal entre mostres consecutives, requisit per a la
        convergència de la xarxa.

  \item \textbf{Xarxa objectiu} $\theta^{-}$: una còpia de $\theta$ que
        s'actualitza periòdicament (cada $N$ passes) i s'usa per calcular els
        objectius de Bellman. Evita l'oscil·lació causada per actualitzar
        simultàniament la xarxa d'estimació i l'objectiu.
\end{itemize}

\begin{figure}[H]
  \centering
  \begin{tikzpicture}[
    bloc/.style={draw, rounded corners=4pt, minimum width=2.5cm,
                 minimum height=0.9cm, align=center, font=\small},
    >=Stealth, thick
  ]
    \node[bloc] (obs)  at (0,   0) {$o_t$};
    \node[bloc] (qnet) at (3.8, 0) {Xarxa Q\\$\theta$};
    \node[bloc] (act)  at (7.6, 0) {$\arg\max_a Q$\\$+\;\mathrm{mask}\to a_t$};
    \node[bloc] (rep)  at (1.9,-2.2) {Replay $\mathcal{D}$};
    \node[bloc, draw=gray!70, text=gray!70] (tnet) at (5.7,-2.2) {$\theta^-$ (objectiu)};

    \draw[->] (obs) -- (qnet);
    \draw[->] (qnet) -- (act);
    \draw[->] (act.south) -- ++(0,-0.55) -| (rep.east);
    \draw[->] (rep.north) |- (qnet.south);
    \draw[->, dashed, gray!70] (qnet.south) |- (tnet.north)
          node[midway, right, font=\scriptsize, text=black]{còpia cada $N$};
    \draw[->, dashed, gray!70] (tnet.east) -| (act.south);
  \end{tikzpicture}
  \caption{Arquitectura \acs{DQN}: la xarxa objectiu $\theta^-$ (discontínua) proporciona objectius de Bellman estables; l'experience replay $\mathcal{D}$ trenca la correlació temporal entre mostres.}
  \label{fig:dqn}
\end{figure}

% ─────────────────────────────────────────────────────────────────────────────
\section{Proximal Policy Optimization}
\label{sec:ppo}

\acf{PPO}~\cite{schulman2017ppo} és un algoritme Actor-Critic on-policy que
s'ha convertit en l'estàndard pràctic de l'\ac{AR} gràcies a la seva
estabilitat i eficiència computacional.

\subsection*{Actor-Critic}

PPO manté dues xarxes: l'\emph{actor} $\pi(a \mid s; \theta)$ que aprèn la
política, i el \emph{crític} $V(s; \phi)$ que estima la funció de valor.
L'actor s'entrena per maximitzar el retorn esperat; el crític proporciona una
línia base que redueix la variança del gradient.

L'avantatge estimat $\hat{A}_t$ mesura si l'acció presa va ser millor o pitjor
del previst pel crític, i s'estima mitjançant el mètode GAE
(\emph{Generalized Advantage Estimation}).

\begin{figure}[H]
  \centering
  \begin{tikzpicture}[
    bloc/.style={draw, rounded corners=4pt, minimum width=2.6cm,
                 minimum height=0.9cm, align=center, font=\small},
    >=Stealth, thick
  ]
    \node[bloc] (obs)    at (0,   0)   {Observació $o_t$};
    \node[bloc] (shared) at (3.8, 0)   {Extractor\\compartit $f_\phi$};
    \node[bloc] (actor)  at (7.6, 0.8) {Actor $\pi_\theta$\\(política)};
    \node[bloc] (critic) at (7.6,-0.8) {Crític $V_\phi$\\(valor)};

    \draw[->] (obs) -- (shared);
    \draw[->] (shared.east) |- (actor.west);
    \draw[->] (shared.east) |- (critic.west);
  \end{tikzpicture}
  \caption{Arquitectura Actor-Critic de \acs{PPO}: extractor de característiques compartit amb dos caps independents, un per a la política i un per a la funció de valor.}
  \label{fig:actor-critic}
\end{figure}

\subsection*{Objectiu amb retall}

La innovació central de PPO és l'objectiu amb retall (\emph{clipped
surrogate objective}), que limita l'actualització de la política per evitar
passos massa grans:
\[
  \mathcal{L}^{\mathrm{CLIP}}(\theta) =
  \mathbb{E}_t\!\left[
    \min\!\left(
      r_t(\theta)\,\hat{A}_t,\;
      \operatorname{clip}(r_t(\theta),\,1-\varepsilon,\,1+\varepsilon)\,\hat{A}_t
    \right)
  \right],
\]
on $r_t(\theta) = \pi(a_t \mid s_t; \theta) / \pi(a_t \mid s_t; \theta_{\mathrm{old}})$
és el ràtio entre la política nova i l'antiga, i $\varepsilon = 0{,}2$.

La pèrdua total afegeix un terme de valor $\mathcal{L}^{V}$ i un bonus
d'entropia $\mathcal{H}[\pi]$ per mantenir l'exploració:
\[
  \mathcal{L} = \mathcal{L}^{\mathrm{CLIP}} - c_1\,\mathcal{L}^{V}
  + c_2\,\mathcal{H}[\pi].
\]

\subsection*{Paral·lelisme}

Com que PPO és \emph{on-policy}, necessita generar mostres noves en cada
iteració. Per compensar aquest cost, l'algorisme aprofita el paral·lelisme
de manera natural: múltiples còpies de l'entorn s'executen simultàniament i
alimenten un únic agent central, tal com proposen Schulman et
al.~\cite{schulman2017ppo}. Aquesta arquitectura permet escalar el nombre de
passos per unitat de temps de forma gairebé lineal amb el nombre d'entorns
paral·lels, compensant l'eficiència de mostres inferior respecte als mètodes
\emph{off-policy}.

% ─────────────────────────────────────────────────────────────────────────────
\section{Curriculum Learning}
\label{sec:curriculum}

El \emph{curriculum learning}~\cite{bengio2009curriculum} és una estratègia
d'entrenament que presenta els exemples o les tasques en ordre creixent de
dificultat, imitant la manera com els humans aprenen.

En el context de l'\ac{AR}, el problema que motiva l'ús del curriculum és el
de la \textbf{recompensa escassa}: en una partida sencera de Truc, l'episodi
pot durar fins a $\sim$150 passos i el senyal de recompensa rellevant arriba
tard i de forma difusa. Amb un horitzó tan llarg, l'agent té dificultats per
atribuir les recompenses a les decisions que les van causar.

\subsection*{Aplicació al Truc}

S'ha dissenyat un currículum de dues etapes:

\begin{enumerate}
  \item \textbf{Etapa de mans} (12M passos): l'agent s'entrena en
        \texttt{TrucGameMa}, on cada mà és un episodi independent de
        $\sim$15 passos amb recompensa densa normalitzada $\text{pts}/24$.
        La recompensa és proporcional als punts guanyats a cada mà.

  \item \textbf{Etapa de partides} (12M passos addicionals): els pesos apresos
        es transfereixen a \texttt{TrucGame}, on l'episodi és una partida
        sencera fins a 24 punts. L'agent parteix d'una política ja formada i
        l'afinament es produeix sobre el comportament estratègic a llarg termini.
\end{enumerate}

\begin{figure}[H]
  \centering
  \begin{tikzpicture}[
    bloc/.style={draw, rounded corners=4pt, minimum width=3.4cm,
                 minimum height=1.3cm, align=center, font=\small},
    >=Stealth, thick
  ]
    \node[bloc] (fase1) at (0,0) {\textbf{Etapa 1: Mans}\\12M passos\\recompensa densa};
    \node[bloc] (fase2) at (6,0) {\textbf{Etapa 2: Partides}\\12M passos\\recompensa escassa};
    \draw[->] (fase1) -- node[above, font=\small]{transferència $\theta$} (fase2);
  \end{tikzpicture}
  \caption{Currículum de dues etapes: mans individuals amb recompensa densa, seguides de partides senceres amb transferència dels pesos apresos.}\label{fig:curriculum}
\end{figure}

El risc principal és l'\textbf{oblit catastròfic}: els algorismes on-policy com
PPO descarten les dades antigues, cosa que pot fer que l'agent oblidi el que va
aprendre a l'etapa de mans quan s'exposa a les partides senceres. 

% ─────────────────────────────────────────────────────────────────────────────
\section{Deep Recurrent Q-Learning}
\label{sec:drqn}

\acf{DQN} assumeix que l'observació $o_t$ és una representació suficient de
l'estat per prendre decisions òptimes. En un \ac{POMDP} això no és cert: sense
memòria de les observacions passades, l'agent pot prendre decisions subòptimes
quan l'estat ocult és rellevant.

\acf{DRQN}~\cite{hausknecht2015drqn} proposa substituir la primera capa densa
de \ac{DQN} per una \acf{GRU} o \acs{LSTM} que manté un estat ocult
$(h_t, c_t)$ propagat a través dels passos de temps. Així, l'aproximació de la
funció de valor depèn no sols de $o_t$ sinó de tota la seqüència
d'observacions passades, el que equival a aproximar la funció de valor sobre la
\emph{creença} (distribució de probabilitat sobre estats ocults).

\subsection*{Entrenament seqüencial}

L'entrenament de \ac{DRQN} propaga els gradients a través del temps (BPTT,
\emph{Backpropagation Through Time}) sobre segments de trajectòria de longitud
$L$. La pèrdua s'acumula als últims $L$ passos de la seqüència per evitar que
els gradients s'extingueixin en contextos molt llargs.

% ─────────────────────────────────────────────────────────────────────────────
\section{Self-play i equilibri de Nash}
\label{sec:selfplay}

En un joc de suma zero entre dos jugadors, un \textbf{equilibri de Nash} és un
parell de polítiques $(\pi^*, \sigma^*)$ tal que cap dels dos jugadors pot
incrementar el seu retorn esperat canviant unilateralment la seva política:
\[
  V(\pi^*, \sigma^*) \geq V(\pi, \sigma^*) \quad \forall\,\pi, \qquad
  V(\pi^*, \sigma^*) \leq V(\pi^*, \sigma) \quad \forall\,\sigma.
\]
Una política en equilibri de Nash és \emph{no explotable}: cap estratègia
adversària pot millorar contra ella.

\subsection*{Fictitious Play i self-play modern}

El \emph{Fictitious Play} (Brown, 1951) és un algorisme iteratiu on cada
jugador best-responds a la política promig de l'adversari. Convergeix a
l'equilibri de Nash en jocs de suma zero de dos jugadors.

El self-play modern~\cite{brown2019pluribus} opera de manera similar: l'agent
s'entrena jugant contra còpies antigues de si mateix (pool de snapshots), la
qual cosa evita l'especialització excessiva contra un únic oponent i afavoreix
la convergència cap a polítiques robustes.

\subsection*{Mètriques de robustesa}

Per mesurar la qualitat de la política més enllà del win rate brut, s'han
definit dues mètriques:
\begin{itemize}
  \item $\texttt{metric\_robust} = \overline{WR} - \lambda\,\sigma(WR_{\text{variant}})$,
        que penalitza la variança del win rate entre variants d'oponent (amb
        $\lambda = 0{,}5$).
  \item $\texttt{exploit\_selfplay} = |WR_{\text{vs snapshots}} - 50|$, que mesura
        la distància a l'equilibri de Nash: 0 indica que cap dels dos jugadors
        pot guanyar sistemàticament a l'altre.
\end{itemize}

% ─────────────────────────────────────────────────────────────────────────────
\section{Neural Fictitious Self-Play}
\label{sec:nfsp}

\acf{NFSP}~\cite{heinrich2016nfsp} és un marc d'entrenament per a jocs
d'informació imperfecta que combina l'\ac{AR} per aprendre la millor resposta
(\emph{best response}) amb aprenentatge supervisat per aprendre la política
promig, aproximant així l'equilibri de Nash sense resoldre explícitament l'arbre
de joc.

\subsection*{Dos bucles en paral·lel}

\begin{enumerate}
  \item \textbf{Bucle RL (millor resposta)}: un agent PPO maximitza la
        recompensa jugant contra el pool d'oponents (que inclou còpies antigues
        i la política promig). Aquest agent aprèn una política que
        \emph{best-responds} a les polítiques de l'historial.

  \item \textbf{Bucle SL (política promig)}: una xarxa auxiliar
        \texttt{AveragePolicyNet} s'entrena per cross-entropy per imitar les
        accions preses pel bucle RL al llarg de tot l'entrenament. Rep com a
        etiquetes les accions del bucle RL i aprèn la distribució promig
        d'accions.
\end{enumerate}

\begin{figure}[H]
  \centering
  \begin{tikzpicture}[
    bloc/.style={draw, rounded corners=4pt, minimum width=3.0cm,
                 minimum height=1.0cm, align=center, font=\small},
    >=Stealth, thick
  ]
    \node[bloc] (rl)  at (0,  1.4) {Bucle RL\\$\pi^{BR}$ (best response)};
    \node[bloc] (buf) at (5.5, 0)  {Reservoir\\Buffer};
    \node[bloc] (sl)  at (0, -1.4) {Bucle SL\\$\overline{\pi}$ (pol. promig)};

    \draw[->] (rl.east)  -- node[above, sloped, font=\scriptsize]{accions de $\pi^{BR}$}
              (buf.north west);
    \draw[->] (buf.south west) -- node[below, sloped, font=\scriptsize]{mostreig uniforme}
              (sl.east);
    \draw[->] (sl.west) -- ++(-1.4,0) |-
              node[left, font=\scriptsize, pos=0.25]{$\overline{\pi}$ com a oponent}
              (rl.west);
  \end{tikzpicture}
  \caption{\acs{NFSP}: el bucle RL genera accions que s'acumulen al reservoir buffer; el bucle SL aprèn per imitació la política promig $\overline{\pi}$, que s'usa com a oponent durant l'entrenament.}\label{fig:nfsp}
\end{figure}

\subsection*{Reservoir Buffer}

Per garantir que l'estimació de la política promig sigui imparcial, les accions
del bucle RL s'emmagatzemen en un \emph{Reservoir Buffer} que aplica mostreig
uniforme de l'historial: cada acció passada té la mateixa probabilitat de
romandre al buffer, independentment del moment en què es va generar.

\subsection*{Paràmetre $\eta$ i Nash gap}

El paràmetre $\eta$ controla la probabilitat que l'oponent jugui la política
promig (SL) en lloc del pool de self-play durant l'entrenament. La distància a
l'equilibri de Nash es mesura com:
\[
  \texttt{exploit\_vs\_sl} = |WR_{\text{PPO vs SL}} - 50|.
\]
Un valor de 0 indica que el best-response i la política promig estan en
equilibri: cap dels dos pot guanyar sistemàticament a l'altre.

\begin{table}[H]
  \centering
  \caption{Comparativa dels algorismes d'\acs{AR} emprats al treball.}\label{tab:comparativa-algorismes}
  \begin{tabular}{lcccc}
    \toprule
    \textbf{Algorisme} & \textbf{On/Off-policy} & \textbf{Memòria} & \textbf{Garantia Nash} & \\
    \midrule
    \acs{DQN}  & Off-policy & Experience replay  & No         \\
    \acs{PPO}  & On-policy  & Cap                & No                 \\
    \acs{NFSP} & Tots dos   & Reservoir buffer   & Aproximada             \\
    \bottomrule
  \end{tabular}
\end{table}
